\documentclass[times, 10pt,twocolumn]{article}
\usepackage{latex8}
\usepackage{times}
\usepackage[utf8]{inputenc}
\usepackage{url}

\pagestyle{empty}

\begin{document}

\title{dida-trade: State-Machine Replication with\\Multi-Paxos, Fast Paxos and Vertical Paxos}

\author{Diogo Noronha da Cruz, n.\textordmasculine~106591 \and Jo\~ao Pedro Correia e Sousa, n.\textordmasculine~99991 \and Tom\'as Dias Monteiro, n.\textordmasculine~106211\\
Instituto Superior T\'ecnico, Universidade de Lisboa\\
Design and Implementation of Distributed Applications 2026/27 --- Group 38
}

\maketitle
\thispagestyle{empty}

%-------------------------------------------------------------------------
\Section{Steps implemented}

We implemented Steps 1--5 (Ultimate goal): the real phase 1 and the debug
modes, Multi-Paxos, Fast Paxos, a causally ordered \texttt{balance}, and
dynamic membership with state transfer (Schedule B with five acceptors, and
Schedule C). The protocol uses no timeouts. Progress comes only from protocol
events, plus manual reconfiguration from the console.

%-------------------------------------------------------------------------
\Section{Multi-Paxos (Steps 1--2)}

A leader runs one phase 1 per ballot, starting at its first unexecuted
position. Acceptors reply with every vote they cast from that position on,
and flag the values they know are decided. For each position the leader
re-proposes the decided value, or else the value voted in the highest ballot.
Gaps are filled with no-ops. Each pending command then gets its own proposer
thread, which reserves the next position under a lock and sends phase 2
without waiting for earlier positions to be decided. Acceptors send their
vote directly to every learner, and positions are executed strictly in order.

\textbf{Auxiliary structures.} \texttt{RequestHistory} keeps client requests
in two maps, \emph{pending} and \emph{processed}, keyed by request id. A
request is added to \emph{pending} when a client sends it, and moved to
\emph{processed} when its position is executed. If the same request id is
decided in a second position, it is found in \emph{processed} and skipped. The
proposer keeps \texttt{proposed\_requests}, the set of request ids it has
already sent to the acceptors. An id is added when a proposer thread is
started for it and removed when it executes. The set is cleared on a ballot
change, so that the new leader proposes again whatever its phase 1 did not
recover. Each log entry (\texttt{PaxosInstance}) keeps the acceptor's own
vote separately from the learner's tally, a map from value to the
\emph{set of acceptors} that voted it. This matters in fast ballots, where
acceptors may vote for different values in the same position.

%-------------------------------------------------------------------------
\Section{Fast Paxos (Step 3)}

The leader of a fast ballot runs phase 1 and then sends an \emph{any} message,
which opens every position from the first free one. Since clients send each
request to every replica, every replica acts as a fast proposer once it has
received the \emph{any}.

\textbf{The acceptor chooses the position.} A fast proposal carries no
position, because without a leader there is no one to serialize reservations.
Each acceptor places the request in its lowest open position. A request it
has already voted for in this ballot keeps its position, and a request it
already knows is decided gets none. Otherwise, a late duplicate would occupy
the next position on that acceptor alone and collide with the following
command. A proposal that reaches an acceptor before the \emph{any} waits for
it, instead of being refused.

\textbf{Quorums.} We extended the \texttt{Schedule} interface with
\texttt{fastquorum} (the super-majority) and \texttt{recoveryquorum} (the
promises needed in phase 1), alongside the existing \texttt{quorum}. With 3
acceptors they are 3 and 2 (classic quorum 2). With 5 acceptors they are 4 and
3 (classic quorum 3).

\textbf{Collisions and recovery.} A learner declares a collision when a fast
quorum of distinct acceptors has voted in a position without agreeing on a
value. It then moves to the next ballot, whose scheduled leader starts phase 1
on its own; this is the automatic reconfiguration of the specification. For
each position, the new leader looks at the highest ballot $k$ voted by its
recovery quorum $Q$. If $k$ was fast, a value $v$ may have been chosen only if
at least $|Q|+F-n$ acceptors of $Q$ voted $v$ in $k$. Such a value must be
re-proposed; otherwise any reported value is safe.

\textbf{Auxiliary structures.} Each proposer keeps \texttt{fast\_marks}, a map
from position to the request it saw voted there in the current fast ballot.
An entry is added when an acceptor reply or a local vote places a request,
and removed when that position is decided. If a different request won the
position, the losing request is also removed from
\texttt{proposed\_requests}, so that it is proposed again.

%-------------------------------------------------------------------------
\Section{Causal \texttt{balance} (Step 4)}

All writes go through one totally ordered log, so a client's causal past fits
in a single integer: the number of log positions it has observed. Each write
reply carries the number of positions the replica had executed. The client
keeps the maximum, \texttt{seen\_slot}, and sends it with \texttt{balance} as
\texttt{min\_slot}. A replica answers \texttt{balance} locally, without
consensus, once it has executed at least \texttt{min\_slot} positions; until
then the request waits. The reply carries the replica's position count, which
the client folds into \texttt{seen\_slot}. The state machine and the position
counter are updated under the same lock that \texttt{balance} reads them with.

%-------------------------------------------------------------------------
\Section{Dynamic membership and state transfer (Step 5)}

\textbf{Completed ballot.} A leader's phase 1 queries the acceptors of every
ballot from the last \emph{complete} ballot $c$ up to its own, and needs a
recovery quorum of each set. A ballot $b$ becomes complete only when three
conditions hold: its leader finished phase 1; every position phase 1
recovered (decided ones included) is accepted by a quorum of $b$'s acceptors;
and, if the acceptor set changed, a quorum of the new acceptors has received
the decided log (\texttt{installstate}). A \texttt{CompletionTracker} per
ballot counts the outstanding recovered positions and the transfer, and only
sets $c$ when both are done. The leader returns $c$ to the console, which
passes it to the next leader with \texttt{ballot} or \texttt{activate}. A
replica that never learns $c$ keeps a lower value, which only makes its phase
1 ask more acceptors.

\textbf{Pull transfer.} Execution blocks at a position that is not decided
while later ones are, or at a position that is decided but whose client
request never reached this replica. The replica then asks every peer
(\texttt{fetchstate}) for the decided positions from there on, \emph{with the
commands}. \texttt{catch\_up\_requested\_instance} records the position it
last asked about, so it asks at most once per blocked position.
\texttt{max\_decided} holds the highest decided position seen, which is what
reveals a gap. We transfer the log rather than a snapshot. Replay is
deterministic and also rebuilds \texttt{RequestHistory}, which a client
request arriving after a transfer needs in order to find its stored reply.

\textbf{Recovery.} A restarted replica has lost its promises and votes. Before
acting as an acceptor, it fetches the log from a recovery quorum of peers and
replays it. If the system was already running, it also abstains from every
ballot up to the highest ballot its peers report (\texttt{abstain\_upto}). It
may have voted in those ballots before crashing, and voting again could
contradict an earlier vote. Abstaining is answered as ``unavailable'', so the
other acceptors can still form a quorum. A peer that receives
\texttt{fetchstate} resets its gRPC reconnect back-off towards the sender,
which could otherwise reject calls to the restarted replica for up to two
minutes.

\textbf{Schedule fixes.} \texttt{ScheduleC} used a classic quorum of 2 out of
5 acceptors, which allowed disjoint quorums; it is now 3. \texttt{ScheduleB}
has five acceptors.

%-------------------------------------------------------------------------
\Section{Limitations}

Without timeouts, a fast ballot that stalls with an acceptor down, or a ballot
whose leader has crashed, is unblocked only by a collision or by a
\texttt{ballot} command from the console. Without stable storage, recovery
relies on abstention. This assumes that some peer has seen every ballot the
crashed replica took part in, which holds because every ballot needs a quorum,
but it is not a proof. The transferred log grows without bound.

%-------------------------------------------------------------------------
\section*{Use of AI}

% TODO(group): this declaration must be accurate. Review it before submitting,
% and state precisely which tools were used and for what.
``All design choices and the code have been by us. The following tools
[list the tools] have been used for this purpose [list the purpose].''

\section*{Authorship}

{\small
We, Diogo Noronha da Cruz, n.\textordmasculine~106591, MEIC; Jo\~ao Pedro
Correia e Sousa, n.\textordmasculine~99991, MEIC; Tom\'as Dias Monteiro,
n.\textordmasculine~106211, MEIC, hereby declare, on our honor, that the
work entitled ``dida-trade'', carried out as part of the Design and
Implementation of Distributed Applications course, resulted from the active
participation of the following group members. The individual contributions
were:
Diogo Cruz --- [summary of contribution];
Jo\~ao Sousa --- [summary of contribution];
Tom\'as Monteiro --- [summary of contribution].
We declare that this work is original and complies with the applicable
academic regulations.

Lisboa, [date] \hfill Signatures:
}

\end{document}
